\begin{summarybox}
\textbf{一句话：} $F\notin\ell$ 时，焦准等距或距离比为 $1$ 刻画抛物线。
\[
PF=PH>0,\qquad e=1,\qquad p=d(F,\ell)>0.
\]
顶点原点、右开口：$y^2=2px$，$F(p/2,0)$，$\ell:x=-p/2$。
$PF=x+p/2\ge p/2$；最短焦半径在顶点取得。
几何上也可识别为“过焦点、切准线的圆”的圆心轨迹。

\textbf{最后四查。}
\begin{enumerate}[leftmargin=*]
\item 焦点是否不在准线上，点到准线是否量垂距？
\item 顶点是否原点；若不是，是否使用平移后的方程？
\item $p$、$p/2$、$2p$ 的几何意义是否一致？
\item 给定焦半径是否至少为 $p/2$，所得点能否代回？
\end{enumerate}
\end{summarybox}
